\documentclass[12pt]{article}
\usepackage{amscd,amssymb,amsthm,amsmath,graphicx, url, color,
verbatim}
\usepackage{enumerate}
\usepackage{parskip}
\usepackage{graphicx} % Required for inserting images

\title{Proof portfolio Math 314}
\author{Yicheng (Eason) Jiang}
\date{August 2026}

\begin{document}

\maketitle


\setlength{\textwidth}{6.5truein}
\setlength{\hoffset}{-.4truein}
\setlength{\textheight}{9.5truein}\setlength\headheight{0in}
\setlength\topmargin{-.75in}
\renewcommand{\(}{\left(}
\renewcommand{\)}{\right)}
\newcommand{\Q}{\mathbb{Q}}
\newcommand{\R}{\mathbb{R}}
\newcommand{\N}{\mathbb{N}}
\newcommand{\Z}{\mathbb{Z}}
\newtheorem{Theorem}{Theorem}
\newtheorem{Conjecture}[equation]{Conjecture}
\newtheorem{Lemma}[equation]{Lemma}
\newtheorem{Proposition}{Proposition}
\newtheorem{Corollary}{Corollary}
\newtheorem*{Definition}{Definition}
\pagestyle{empty}
\begin{center}
{\bf Proof portfolio part I}
\end{center}
\begin{Lemma} If $n$ is any integer, then $n^2+n$ is even.
\end{Lemma}
\begin{proof}
Let $n \in \Z$.

We have $n^2+n=n(n+1)$.

Assume $n$ is even, so $\exists k_1 \in \Z, \text{s.t. } n=2k_1$.
\[
n^2+n=2k_1(n+1)
\]
$\exists k_2 = k_1(n+1) \in \Z, \text{s.t. } n^2+n =2k_2$, so $n^2+n$ is even if $n$ is even.

Assume $n$ is odd, so $\exists k_3 \in \Z, \text{s.t. } n=2k_3+1$.
\begin{align*}
n+1 &= 2k_3+1+1=2(k_3+1) \\
n^2+n &= n(n+1)=n(2k_4)=2(nk_4)
\end{align*}
$\exists k_4=k_3+1 \in \Z, \text{s.t. } n+1=2k_4$, so $n+1$ is even.

$\exists k_5=nk_4\in \Z, \text{s.t. } n^2+n = 2k_5$.

So $n^2+n$ is even if $n$ is odd.

Overall $n^2+n$ is even if $n \in \Z$.
\end{proof}
\vfill
\begin{Lemma} Let $x$ and $y$ be integers. Prove that $x^{2}+x
y$ is odd if and only if $x$ is odd and $y$ is even.
\end{Lemma}
\begin{proof}
(Forward) (Proof by contradiction) Let $x^2+xy$ is odd.

$\exists k_1 \in \Z, s.t. x^2+xy = 2k_1+1$

If x is even:
$\exists k_2 \in \Z, s.t. x = 2k_2$

\[
x^2+xy = 4 k_2^2+2k_2y = 2(2k_2^2+k_2y)
\]

So, $\exists k_3=2k_2^2+k_2y \in \Z, s.t. x^2+xy = 2 k_3$, $x^2+xy$ is even, contridiction, x is odd.

If y is odd:
$\exists k_4 \in \Z, s.t. y = 2k_4 +1$

$x^2+xy= x^2+(2k_4+1)x = x(x+2k_4+1)$

Since x is odd (Proof above): 
$\exists k_5 \in \Z, s.t. x = 2k_5 +1$

$x(x+2k_4+1)=(2k_5+1)(2k_5+2k_4+2)=2[(2k_5+1)(k_5+k_4+1)]$

$\exists k_6= (2k_5+1)(k_5+k_4+1) \in \Z, s.t. x(x+2k_4+1) = 2k_6$

So $x(x+2k_4+1)$ is even, contradiction.

So $x^2 + xy$ is odd if $x$ is odd and $y$ is even. 

(Backward) Let x be odd and y be even.

$\exists k_7 \in \Z, s.t. x = 2k_7+1$

$\exists k_8 \in \Z, s.t. y = 2k_8$
\[
x^2+xy = x(x+y)=(2k_7+1)(2k_7+1+2k_8)
\]
\[
=4k_7^2+2k_7+2k_7+1+4k_7k_8+2k_8
\]
\[
=2(2k_7^2+k_7+k_7+2k_7k_8+k_8)+1
\]

$\exists k_9= 2k_7^2+2k_7+2k_7k_8+k_8\in \Z, s.t. x^2+xy=2k_9+1$

So $x^2+xy$ is odd. 
\end{proof}
\vfill
\begin{Lemma} If $x$ and $y$ are real numbers with $x y>0$,
then $\frac{x}{y}+\frac{y}{x} \geq 2$. Determine when equality
holds.
\end{Lemma}
\begin{proof}
Since $(x-y)^2 \geq 0$,

\[
\Rightarrow x^2-2xy+y^2 \geq 0
\]

\[
\Rightarrow x^2+y^2 \geq 2xy
\]

Since $xy>0$,

\[
\frac{x^2+y^2}{xy} \geq \frac{2xy}{xy}
\]

\[
\Rightarrow \frac{x}{y}+ \frac{y}{x}\geq 2
\]

The equality holds when $(x-y)^2=0$, $x=y$.

\end{proof}
\vfill
\hrulefill
\vfill
Let $(x_n) = (x_1,x_2,x_3,...)$ denote a sequence of real
numbers. Consider the following definitions\\
\begin{Definition}
"$(x_n)$ diverges to $\infty$" means:
$ \forall M > 0, \exists N \in \mathbb{R}$ such that $n > N
\Rightarrow x_n > M$
\end{Definition}
\begin{Definition}
"$(x_n)$ converges to $L$" means: $\forall \epsilon > 0$, $\exists N \in \mathbb{R}$ such
that $n > N \Rightarrow |x_n - L| < \epsilon$
\end{Definition}
\begin{Definition}
"$(x_n)$ does not to diverge to $\infty$" means: $ \exists M > 0, \forall N \in \mathbb{R}$ such that $n > N$
and $x_n \leq M$
\end{Definition}
\begin{Definition}
"$(x_n)$ does not converge to $L$" means: $\exists \epsilon > 0$, $\forall N \in \mathbb{R}$ such
that $n > N$ and $|x_n - L| \geq \epsilon$
\end{Definition}
\begin{Definition}
"$(x_n)$ does not to converges at all" means: $\forall L \in \R$, $\exists \epsilon > 0$, $\forall N \in \mathbb{R}$ such
that $n > N$ and $|x_n - L| \geq \epsilon$
\end{Definition}
\vfill
\hrulefill
\vfill
In the next two Lemmas, all variables are integers.
\begin{Lemma} If $a\mid b$ and $c\mid d$ then $ac\mid bd$.
\end{Lemma}
\begin{proof}
$\exists k_1 \in \Z$, such that $\frac{b}{a} = k_1$, and

$\exists k_2 \in \Z$, such that $\frac{d}{c} = k_2$.

So $k_1k_2 = \frac{bd}{ac}$,

and $k_1k_2 \in \Z$.

So $\exists k_3=k_1k_2 \in \Z$, such that $\frac{bd}{ac} = k_3$.

$\Rightarrow ac \mid bd$. 
\end{proof}
\vfill
\begin{Lemma}
If $a\mid b$ and $a\mid c$ then $a\mid (bx+cy)$ for any $x, y\in
\Z$.
\end{Lemma}
\begin{proof}
$\exists k_1 \in \Z$, such that $\frac{b}{a} = k_1$, and

$\exists k_2 \in \Z$, such that $\frac{c}{a} = k_2$.

$\forall x,y \in \Z$:

\[
xk_1+yk_2 \in \Z
\]
\[
x(\frac{b}{a})+y(\frac{c}{a})=\frac{bx+cy}{a} \in \Z
\]

So $a\mid (bx+cy)$ for any $x, y\in
\Z$.
\end{proof}
\vfill
\pagestyle{empty}
\begin{center}
{\bf Proof portfolio part II}
\end{center}
\begin{Lemma} $4^{n}-1$ is divisible by 3 for every natural number $n$.
\end{Lemma}
\begin{proof}
(By Induction) 

1. Base cases: 

$4^1-1 = 3$, $3 \mid 3$.

2. Induction Steps:
Let $n \in \N$, $4^n-1$ is divisible by 3.

$\Rightarrow 4^n-1 \equiv 0$ (mod 3)

$\Rightarrow 4^n \equiv 1$ (mod 3)

$4^{n+1}-1 \equiv 4(1)-1 \equiv 0$ (mod 3)

$\Rightarrow 4^{n+1}-1$ is divisible by 3.

$4^{n}-1$ is divisible by 3, $\forall n \in \N$.
\end{proof}
\vfill
\begin{Definition}
Let $A=\left\{(x, y) \in \mathbb{R}^{2}: 2-x-2 y<(2-x)(1-y)\right\}$.

$2-x-2y< 2-x-2y+xy$
$0<xy \Rightarrow xy>0$

So x, y have the same sign.

\end{Definition}
\vfill
\begin{Lemma} $$
(A \cup B)-C \subseteq(A-(B \cup C)) \cup(B-(A \cap C))
$$
\end{Lemma}
\begin{proof}
Let $x \in$ Left Hand Side.
$x \in (A \cup B)-C \iff ((x \in A) \vee (x \in B)) \wedge (x \notin C)$ 

if $x \in A$:

$x\notin B$, by $(x \in A) \vee (x \in B)$.
Since $x \notin C$, $x\notin (B \cup C)$

$\Rightarrow x\in A - (B \cup C)$

if $x \in B$:

$x\notin A$.

$\Rightarrow x \in B - (A \cap C)$
 
So $x \in (A-(B \cup C)) \vee x \in (B-(A \cap C))$

$\Rightarrow x \in (A-(B \cup C)) \cup(B-(A \cap C))$

So, $(A \cup B)-C \subseteq(A-(B \cup C)) \cup(B-(A \cap C))$.

\end{proof}
\vfill
\begin{Lemma} Let $f : A \rightarrow B$ be a function and $X,Y \subseteq A$.
Then $f(X \bigcap Y) \subseteq f(X) \bigcap f(Y)$ with equality when $f$ is injective.
\end{Lemma}
\begin{proof}
Let $b \in f(X \bigcap Y)$, 

then $\exists a\in X \cap Y$, such that $f(a)=b$

$\Rightarrow (a \in X) \wedge (a\in Y)$

$b = f(a) \in f(X)$ and $b = f(a) \in f(Y)$

$b\in f(X) \bigcap f(Y)$

$\Rightarrow f(X \bigcap Y) \subseteq f(X) \bigcap f(Y)$

If f is injective: 

Let $d \in f(X) \bigcap f(Y)$

$\exists c_1 \in X, c_2 \in Y, f(c_1) = f(c_2) = d$

Since f is injective, $c_1 = c_2$.

$\Rightarrow c_1 \in X$ and $c_1 = c_2 \in Y$

So $c_1 \in X \cap Y$

Where $d = f(c_1) \in f(X \bigcap Y)$

So $f(X) \bigcap f(Y) \subseteq f(X \bigcap Y)$.

$f(X) \bigcap f(Y) = f(X \bigcap Y)$
\end{proof}
\vfill
\begin{Lemma}
Suppose that $q \neq 1$ is a real number. Then
$$
1+q+q^{2}+\cdots+q^{n}=\frac{1-q^{n+1}}{1-q} \quad \text { for all } n \in
\mathbb{Z}, n \geq 0
$$
\end{Lemma}
\begin{proof}
(Induction)

Base cases:
\[
1 + q^0 = \frac{1-q^{1}}{1-q}
\]
$1=1$, true.

Induction steps:

Assume: $$
1+q+q^{2}+\cdots+q^{n}=\frac{1-q^{n+1}}{1-q} \quad$$

$$
1+q+q^{2}+\cdots+q^{n} + q^{n+1}=\frac{1-q^{n+1}}{1-q} +q^{n+1}\quad$$

$$
1+q+q^{2}+\cdots+q^{n} + q^{n+1}=\frac{1-q^{n+1}}{1-q} +\frac{q^{n+1} (1-q)}{1-q}\quad$$

$$
1+q+q^{2}+\cdots+q^{n} + q^{n+1}=\frac{1-q^{n+1}+q^{n+1}-q^{n+2}}{1-q}\quad$$

$$
1+q+q^{2}+\cdots+q^{n} + q^{n+1}=\frac{1-q^{n+2}}{1-q}\quad$$

So $$
1+q+q^{2}+\cdots+q^{n}=\frac{1-q^{n+1}}{1-q} \quad \text { for all } n \in
\mathbb{Z}, n \geq 0
$$
\end{proof}
\vfill
\begin{Lemma}
Let $f$ be a function such that $f(x y)=x f(y)+y f(x)$ for all $x, y \in \mathbb{R}
$. Then $f(1)=0$ and $f\left(u^{n}\right)=n u^{n-1} f(u)$ for all $u \in \mathbb{R}
$ and $n \in \mathbb{N}$.
\end{Lemma}
\begin{proof}
\[
f(1) = f(1\cdot1) = 1f(1)+1f(1) = 2f(1)
\]

\[
2f(1) - f(1) = f(1) = 0
\]
(Induction)

Base cases: 

\[
f(u^1) = f(u)
\]

\[
1u^{0}f(u) = f(u)
\]

Equality holds.

Induction steps:

Assume $f\left(u^{n}\right)=n u^{n-1} f(u)$.

\[
f(u^{n+1}) = f(u\cdot u^n) = uf(u^n)+u^nf(u) = u \cdot nu^{n-1}f(u)+u^nf(u)
\]

\[
=(nu^n+u^n)f(u) = (n+1)u^nf(u)
\]

So, $f\left(u^{n}\right)=n u^{n-1} f(u)$ for all $u \in \mathbb{R}
$ and $n \in \mathbb{N}$.
\end{proof}
\vfill
\begin{Lemma}
Let $P$ be the set of polynomials with real coefficients and with degree $\geq
1$. If $f \in P$,let $A_{f}:=\{x \in \mathbb{R}: f(x)=0\}$.\\
Then $\bigcup\limits_{f \in P} A_{f}=\mathbb{R}$ and $\bigcap\limits_{f \in P}
A_{f}=\emptyset$.
\end{Lemma}
\begin{proof}
I. $\bigcup\limits_{f \in P} A_{f}=\mathbb{R}$

i. $\R \subseteq \bigcup\limits_{f \in P} A_{f}$.

$\forall a \in \R$, 
$f_a(x) := x - a \in P$.

\[
f_a(a) = a-a = 0
\]
So, $\forall a \in \R, \exists f_a \in P, f_a(a) = 0$.

$\Rightarrow \forall a\in \R,a \in \bigcup\limits_{f \in P} A_{f}$

$\R \subseteq \bigcup\limits_{f \in P} A_{f}$.

ii. $\bigcup\limits_{f \in P} A_{f} \subseteq \R $.

Assume $b \in \bigcup\limits_{f \in P} A_{f}$.

$\forall f \in P, A_f \subseteq \R$

$b \in A_f$, for some $f\in P$, so $b \in \R$.

$\bigcup\limits_{f \in P} A_{f} \subseteq \R $

By i. and ii., $\bigcup\limits_{f \in P} A_{f} = \R $.

II. $\bigcap\limits_{f \in P}
A_{f}=\emptyset$.

(Contradiction)

Assume $\exists c \in R, \forall f \in P, f(c) = 0$.

Let $f_1(x) :=x-1 \in P$, $f_2(x) :=x-2 \in P$.

Since $f_1(c) = 0$, $c-1=0$, $c=1$.

Since $f_2(c) = 0$, $c-2=0$, $c=2$.

Contradiction, there is no such $c$ in $\bigcap\limits_{f \in P}
A_{f}$.

$\iff \bigcap\limits_{f \in P}
A_{f} = \emptyset$
\end{proof}
\vfill
\begin{Lemma}
The relation $\sim$ defined on $\Z\times \N$ by
\[(a,b)\sim (c,d)\iff ad = bc\]
is an equivalence relation.
\end{Lemma}
\begin{proof}
i. Reflexive 

Let $a\in \Z$, and $b\in \N$.
\[(a,b)\sim (a,b)\iff ab = ba\]

And $ab = ba$ by properties of multiplication under $\R$, and $\Z,\N \subset \R$, so $\sim$ is reflexive

ii. Symmetric

Let $c\in \Z$, and $d\in \N$, assume
\[(a,b)\sim (c,d)\]
Where $ad = bc$.

Since $ad = bc \iff da = cb$,
\[da = cb\iff (c,d)\sim (a,b)\]
So $\sim$ is symmetric.

iii. Transitive

Let $e\in \Z$, and $f\in \N$, assume $(a,b)\sim (c,d)$ and $(c,d) \sim (e,f)$.

$\Rightarrow ad = bc, cf = de$.

Since $b \in \N$, $b \neq 0$.

$c= \frac{ad}{b}$, and substitute it in second equation.
\[\frac{ad}{b}f = de\]
\[adf = dbe\]
Since $d\in \N$, $d \neq 0$.
\[af = de\]
\[
(a,d) \sim (e,f)
\]

$\sim$ is transitive.

So $\sim$ is an equivalence relation.

\end{proof}
\vfill
\end{document}
